\scanpage{23}
fiber $(\mathfrak X_{\bar s},\underline I_{\bar s})$ is a geometric \inserted{MD-}curve over $k(s)$ -- i.e. $\mathfrak X_{\bar s}$ is $0$-connected, of dimension $1$, \deleted{\unsure{as well as}} the singular points are \inserted{$\notin I_{\bar s}$ and are} ordinary double points, and finally the condition of infinitesimal rigidity \inserted{(in \unsure{automorphisms|conditions})} is satisfied -- \deleted{thus} i.e. one requires that $\mathfrak X_{\bar s}$ not be smooth of genus $1$ with $I_{\bar s}=\unsure{\phi|\varnothing}$, and that the irreducible components of genus $0$ of $\mathfrak X_{\bar s}$ be of ``marked weight''
\[
\geqslant 3.
\]
The difference from the considerations of the preceding number is that now ($\bar m$ over a connected base $\mathcal S$, for example \deleted{$\unsure{\mathcal S=}$} a trait) the \uline{types} of the geometric fibers can vary from one geometric point to another. But it is clear that the \uline{numerical} type $(g,\nu)$ does not change (invariance \unsure{of the|of} $\chi(\widehat X_s,\underline O_{\widehat X_s})\ldots$ and of $\operatorname{card}\underline I_s$), more precisely,
% Two clipped right-edge ink fragments; wording and extent unrecoverable.
% Lower fragment is crossed; its deletion state remains less than certain.
\margin{\illegible{1}\ \deleted{\illegible{1}}}
% Light lower-edge pagination; hand unidentified.
\mspage{22}
